
\chapter{Anhang: Grobablauf für 2 Lektionen (Sek II, 4. Klassen)}

Dieser Ablauf ist als pragmatische Orientierung für das Wahlpflichtmodul Statistik gedacht.
Die zeitliche Planung ist bewusst als Richtwert formuliert und kann je nach Lerngruppe angepasst werden.

\section{Lektion 1 (45 Minuten): Von der Frage zu ersten Befunden}

\begin{table}[H]
	\centering
	\small
	\begin{tabular}{p{0.14\textwidth} p{0.26\textwidth} p{0.12\textwidth} p{0.42\textwidth}}
		\toprule
		\textbf{Zeitfenster} & \textbf{Phase} & \textbf{Lehrperson} & \textbf{Inhalt und Ziel} \\
		\midrule
		0--5 min  & Einstieg & \texttt{sb} & Kurzer Impuls mit einer alltagsnahen Statistikfrage; Ziel ist die Aktivierung von Vorwissen und die Klärung, was eine statistische Frage von einer Einzelfrage unterscheidet. \\
		5--12 min & Fragestellung schärfen & \texttt{sb} & Sammeln und Umformulieren erster Fragen der Lernenden: von Personen- oder Themenfragen hin zu Variablen, Verteilungen oder Vergleichen. \\
		12--20 min & Datensatz einführen & \texttt{wn} & Einführung in den Datensatz (z.B. Schlaf und Fokus), Zielgruppe, Variablen, Skalenniveaus und mögliche Grenzen der Aussagekraft. \\
		20--35 min & Erste Auswertung & \texttt{wn} & Lernende bestimmen in Partnerarbeit zentrale Kennzahlen (z.B. Median, Mittelwert, Minimum, Maximum) und begründen die Wahl passender Kennzahlen. \\
		35--43 min & Visualisierung & \texttt{wn} & Erste grafische Darstellung (Histogramm oder Boxplot) und kurze Interpretation: typische Werte, Streuung, mögliche Ausreisser. \\
		43--45 min & Sicherung & \texttt{wn} & Gemeinsame Kurzsicherung: Was ist bereits belastbar beschreibbar, was bleibt offen? \\
		\bottomrule
	\end{tabular}
	\caption{Grobablauf Lektion 1}
\end{table}

\section{Lektion 2 (45 Minuten): Vergleich, Zusammenhang und begründete Aussage}

\begin{table}[H]
	\centering
	\small
	\begin{tabular}{p{0.14\textwidth} p{0.26\textwidth} p{0.15\textwidth} p{0.42\textwidth}}
		\toprule
		\textbf{Zeitfenster} & \textbf{Phase} & \textbf{Lehrperson} & \textbf{Inhalt und Ziel} \\
		\midrule
		0--6 min  & Rückblick & \texttt{wn} & Kurze Reaktivierung der Ergebnisse aus Lektion 1 und Präzisierung der Anschlussfrage (Vergleich oder Zusammenhang). \\
		6--18 min & Vertiefung Analyse & \texttt{wn} & Gruppenvergleich oder Zusammenhangsanalyse (z.B. Boxplot-Vergleich oder Scatterplot SleepHours vs FocusScore) mit Fokus auf vorsichtige Interpretation. \\
		18--30 min & Schliessende Perspektive & \texttt{sb} & Einfache Hinführung zur Unsicherheit von Aussagen (Stichprobe vs Grundgesamtheit, Idee von Signifikanz bzw. Vertrauensintervall auf konzeptueller Ebene). \\
		30--40 min & Transferauftrag & \texttt{sb} & Lernende formulieren eine eigene statistische Fragestellung mit Zielgruppe, Variablen und geplanter Auswertung. \\
		40--45 min & Abschluss & \texttt{sb} & Kurzpräsentationen ausgewählter Fragestellungen und gemeinsames Feedback zu Klarheit, Messbarkeit und Aussagebereich. \\
		\bottomrule
	\end{tabular}
	\caption{Grobablauf Lektion 2}
\end{table}

\section{Erwartete Produkte nach den 2 Lektionen}

Am Ende der beiden Lektionen liegt idealerweise eine klar formulierte statistische Fragestellung mit definierter Zielgruppe vor,
ergänzt durch erste deskriptive Kennzahlen, mindestens eine passende Visualisierung und eine kurze, fachlich vorsichtige Interpretation.
Zusätzlich ist transparent, welche Aussagen durch die vorliegenden Daten gestützt sind und welche nicht.

\restoregeometry